% Corrected French for printed p.141.
% Corrections: coordinate the two searches and restore duaux; 0<=i<=a in
% Proposition 2.2.1; à partir de; T^0(A') as the domain of f^0(A'); standards.
\SourcePage{141}

\LayerTarget{tohoku:remark:satellites-existence:p2}
\noindent que pour tout foncteur covariant additif $F$ de $C$ dans $C'$, il
existe un $\partial$-foncteur universel défini en tous les degrés et se
réduisant à $F$ pour le degré zéro (i. e. tous les satellites $S^iF$
existent) il suffit manifestement de prouver que $S^1F$ et $S^{-1}F$ existent,
en vertu des formules
\[
 S^1(S^iF)=S^{i+1}F\quad\hbox{si }i\geq0,\qquad
 S^{-1}(S^iF)=S^{i-1}F\quad\hbox{si }i\leq0,
\]
\noindent (dont la vérification à partir de la définition est triviale).
D'ailleurs, la recherche de $S^1F$ et celle de $S^{-1}F$ sont deux problèmes
duaux, ces
deux foncteurs se permutent si on remplace $C$ et $C'$ par les catégories
duales.

\LayerTarget{tohoku:def:effaceable-coeffaceable}
Un foncteur additif $F$ de $C$ dans $C'$ est dit \textit{effaçable} si pour
tout $A\in C$, on peut trouver un monomorphisme $u:A\longrightarrow M$ tel que
$F(u)=0$ ; si $C$ est tel que tout objet $A\in C$ admet un effacement injectif
(cf. 1.10 remarque 1) il revient au même de dire que $F(u)=0$ pour tout
effacement injectif $u$ ; si enfin $C$ est tel que tout objet $A\in C$ admet
un monomorphisme dans un objet injectif $M$ (cf. 1.10.1) il revient au même de
dire que $F(M)=0$ pour tout objet injectif $M$. Dualement, $F$ est dit
\textit{coeffaçable} si pour tout $A\in C$, on peut trouver un épimorphisme
$u:P\longrightarrow A$ tel que $F(u)=0$.

\LayerTarget{tohoku:prop:2.2.1}
\noindent\textsc{Proposition 2.2.1.} --- \textit{Soient $C,C'$ deux catégories
abéliennes, $T=(T^i)_{0\leq i\leq a}$ un $\partial$-foncteur exact (covariant
ou contravariant) de $C$ dans $C'$, avec $a>0$. Si $T$ est covariant et si
$T^i$ est effaçable pour $i>0$, alors $T$ est un $\partial$-foncteur universel,
et la réciproque est vraie si $C$ est tel que tout objet $A\in C$ admet un
effacement injectif (cf. 1.10). Si $T$ est contravariant, le même énoncé vaut
en remplaçant « effaçable » et « effacement injectif » par « coeffaçable » et
« effacement projectif », respectivement.}

\LayerTarget{tohoku:prop:2.2.1:proof:p1}
Il suffit pour la partie directe, de prouver par exemple que si
$(T'^0,T'^1)$ est un $\partial$-foncteur défini pour les degrés 0, 1 et $f^0$
un morphisme fonctoriel $T^0\longrightarrow T'^0$, alors il existe un
morphisme $(f^0,f^1)$ unique de $(T^0,T^1)$ dans $(T'^0,T'^1)$ se réduisant à
$f^0$ pour le degré 0 (nous avons supposé $T$ covariant pour fixer les idées).
Soit $A\in C$, considérons une suite exacte
$0\longrightarrow A\longrightarrow M\longrightarrow A'\longrightarrow 0$
telle que le morphisme $T^1(A)\longrightarrow T^1(M)$ soit nul. Si on a pu
construire $f^1$, on aura un diagramme commutatif

\LayerTarget{tohoku:diagram:prop-2.2.1}
\[
\begin{array}{ccccccc}
T^0(M)&\longrightarrow&T^0(A')&\longrightarrow&T^1(A)&\longrightarrow&T^1(M)\\
\big\downarrow&&\big\downarrow&&\big\downarrow&&\\
T'^0(M)&\longrightarrow&T'^0(A')&\longrightarrow&T'^1(A)&&
\end{array}
\]

\LayerTarget{tohoku:prop:2.2.1:proof:p2}
La première ligne étant exacte, on en conclut que le morphisme
$T^0(A')\longrightarrow T^1(A)$ est surjectif, et par suite que le morphisme
$f^1(A):T^1(A)\longrightarrow T'^1(A)$ est complètement déterminé par passage
au quotient à partir de
$f^0(A'):T^0(A')\longrightarrow T'^0(A')$, ce qui prouve l'unicité de
$f^1(A)$. D'ailleurs, le diagramme précédent sans $f^1(A)$, compte tenu du
fait que le produit des deux morphismes de la deuxième ligne est nul, permet
de définir un morphisme $T^1(A)\longrightarrow T'^1(A)$ de façon unique par la
condition que le diagramme reste commutatif. Des raisonnements standards
montrent que le morphisme ainsi défini ne dépend pas du choix particulier de
la suite exacte
$0\longrightarrow A\longrightarrow M\longrightarrow A'\longrightarrow 0$,
puis le fait que ce morphisme est fonctoriel, et
\textquotedblleft permute à $\partial$\textquotedblright. Cela prouve donc la
première partie de la proposition. La deuxième partie est contenue dans le
théorème d'existence suivant :

\LayerTarget{tohoku:thm:2.2.2:p1}
\noindent\textsc{Théorème 2.2.2.} --- \textit{Soit $C$ une catégorie
abélienne telle que tout objet $A\in C$ admette un effacement injectif
(cf. 1.10). Alors pour tout foncteur additif covariant $F$ sur $C$, les
satellites $S^iF$ ($i\geq0$) existent et sont des foncteurs effaçables pour
$i>0$. Pour que le $\partial$-foncteur universel $(S^iF)_{i\geq0}$ soit exact,
il}

\TohokuContinuation{142}
